\chapter*{Nomenclature}
The list below gives the abbreviations used in this report.

\vspace{0.5cm}
\begin{tabular}{p{3cm} p{11cm}}
BERT & Bidirectional Encoder Representations from Transformers\\
CI & Confidence Interval\\
CPU / GPU & Central Processing Unit / Graphics Processing Unit\\
CUAD & Contract Understanding Atticus Dataset\\
ECE & Expected Calibration Error\\
EM & Exact Match\\
FN / FP & False Negative / False Positive\\
IP & Intellectual Property\\
LLM & Large Language Model\\
LR & Logistic Regression\\
MIL & Multiple-Instance Learning\\
MPS & Metal Performance Shaders (Apple GPU backend)\\
NLP & Natural Language Processing\\
QA & Question Answering\\
ROUGE & Recall-Oriented Understudy for Gisting Evaluation\\
RQ & Research Question\\
SD & Standard Deviation\\
SEC & US Securities and Exchange Commission\\
SQuAD & Stanford Question Answering Dataset\\
TF--IDF & Term Frequency--Inverse Document Frequency\\
TN / TP & True Negative / True Positive\\
\end{tabular}

\clearpage
\textbf{Glossary of terms}
\begin{description}[itemsep=4pt]
  \item[AND-ensemble] A rule that reports a clause category only when both the TF--IDF model and the DistilBERT model say it is present.
  \item[Bootstrap] A way to measure uncertainty: the test contracts are drawn again at random, with repeats, many times, and the score is recomputed each time.
  \item[Calibration] Whether a model's score can be read as a real chance. A well-calibrated score of 0.8 is right about 80\% of the time.
  \item[Clause category] One of the 41 types of contract clause labeled in CUAD (for example \emph{Governing Law}, \emph{Non-Compete}, \emph{Cap on Liability}).
  \item[Contract-level split] A train/test split that puts each whole contract on one side, so no clause from a training contract can appear in the test set.
  \item[Max-pooling (over windows)] Scoring every window of a document and using the highest window score as the score for the whole document.
  \item[Micro-F1] One F1 score computed over all (contract, category) decisions together, so every decision counts the same.
  \item[Micro-F2] The same, but treating recall as twice as important as precision, so a missed clause costs more than a false alarm.
  \item[Presence classification] Deciding, for each clause category, whether a contract contains at least one clause of that category.
  \item[Recall-first setting] A setup whose rule and thresholds were chosen on the validation split to maximize micro-F2; Clausify uses one by default.
  \item[Retrieval ceiling] The most a ``search first'' design can ever find: a clause in a window that the search never picks can never be found.
  \item[Span extraction] Finding the exact text of a clause inside the contract, treated as a question-answering task.
  \item[Token-F1] How much the extracted text overlaps the real clause, word by word (1.0 means a perfect match).
  \item[Validation split] The 81 training contracts held out in the re-run to choose settings such as thresholds, so that the test contracts play no part in any choice.
  \item[Weaker party] The person or company signing the contract with less bargaining power; the risk table is written from their side.
\end{description}
